\section*{Chapter \ref{ch:b1:integration} --- Integration on a Segment}

\begin{solution}{exo:b1:integration:1}
$\dfrac{1}{x^2-4} = \dfrac{1/4}{x - 2} - \dfrac{1/4}{x+2}$
(cover-up), so
\[
\int_0^1 \frac{\dd x}{x^2 - 4}
= \frac14\Bigl[\ln\abs{x-2} - \ln\abs{x+2}\Bigr]_0^1
= \frac14\Bigl(\ln\frac{1}{3} - \ln 1\Bigr) = -\frac{\ln 3}{4}.
\]

By parts ($u' = 1$, $v = \ln t$): $\int_1^{\eu} \ln t\,\dd t =
\bigl[t\ln t\bigr]_1^{\eu} - \int_1^{\eu} \dd t = \eu - (\eu - 1) =
1$.

By parts ($u' = \sin t$, $v = t$): $\int_0^\pi t\sin t\,\dd t =
\bigl[-t\cos t\bigr]_0^\pi + \int_0^\pi \cos t\,\dd t = \pi + 0 =
\pi$.
\end{solution}

\begin{solution}{exo:b1:integration:2}
Substitution $u = t^2 + 1$, $\dd u = 2t\,\dd t$:
\[
\int_0^1 \frac{t\,\dd t}{(t^2+1)^2}
= \frac12 \int_1^2 \frac{\dd u}{u^2}
= \frac12\Bigl[-\frac1u\Bigr]_1^2 = \frac14 .
\]
With $u = \sin t$: $\int_0^{\pi/2} \cos^3 t\,\dd t = \int_0^{\pi/2}
(1 - \sin^2 t)\cos t\,\dd t = \bigl[\sin t - \frac{\sin^3
t}{3}\bigr]_0^{\pi/2} = 1 - \frac13 = \frac23$.
\end{solution}

\begin{solution}{exo:b1:integration:3}
$u_n = \dfrac1n \sum_{k=1}^{n} \dfrac{1}{1 + (k/n)^2}$: a Riemann sum
of $x \mapsto \frac{1}{1+x^2}$ on $\intcc{0}{1}$, so $u_n \to
\int_0^1 \frac{\dd x}{1+x^2} = \arctan 1 = \dfrac\pi4$.

$\ln v_n = \dfrac1n \sum_{k=1}^{n} \ln\dfrac{n+k}{n} = \dfrac 1n
\sum_{k=1}^n \ln\Bigl(1 + \dfrac kn\Bigr) \to \int_0^1 \ln(1+x)\,\dd
x = \bigl[(1+x)\ln(1+x) - x\bigr]_0^1 = 2\ln 2 - 1$. Hence $v_n \to
\eu^{2\ln 2 - 1} = \dfrac 4\eu$. (Check of the identification:
$\frac{(2n)!}{n!\,n^n} = \prod_{k=1}^{n} \frac{n+k}{n}$.)
\end{solution}

\begin{solution}{exo:b1:integration:4}
The limit is $0$. Let $M = \sup \abs f$ and $\varepsilon \in
\intoo{0}{1}$. Cut at $1 - \varepsilon$:
\[
\Bigl| \int_0^1 x^n f \Bigr|
\leq \int_0^{1 - \varepsilon} x^n \abs f + \int_{1-\varepsilon}^1
x^n \abs f
\leq M\,(1-\varepsilon)^n + M\varepsilon .
\]
Since $(1 - \varepsilon)^n \to 0$ (\cref{exo:b1:seq:3}), the limsup
of the left side is $\leq M\varepsilon$ for every $\varepsilon$:
the integral tends to $0$.
\end{solution}

\begin{solution}{exo:b1:integration:5}
$Q(\lambda) = \int_a^b (f + \lambda g)^2 = \int f^2 + 2\lambda \int
fg + \lambda^2 \int g^2 \geq 0$ for all $\lambda$. If $\int g^2 =
0$, then $g = 0$ (strict positivity,
\cref{thm:b1:integration:props}~(4)) and the inequality is $0 \leq
0$. Otherwise $Q$ is a genuine quadratic, everywhere $\geq 0$: its
discriminant is $\leq 0$, i.e.\ $\bigl(\int fg\bigr)^2 \leq \int
f^2 \int g^2$.

Equality iff the discriminant vanishes iff $Q(\lambda_0) = 0$ for
some $\lambda_0$, i.e.\ $\int (f + \lambda_0 g)^2 = 0$, i.e.\ (strict
positivity again) $f = -\lambda_0 g$: equality holds exactly when
$f$ and $g$ are proportional.
\end{solution}

\begin{solution}{exo:b1:integration:6}
Let $\Phi(a) = \int_a^{a+T} f$. By the fundamental theorem
(\cref{thm:b1:integration:ftc}), $\Phi$ is differentiable with
$\Phi'(a) = f(a + T) - f(a) = 0$: constant.

For $x > 0$, write $x = nT + r$, $0 \leq r < T$ ($n = \lfloor x/T
\rfloor$). Chasles:
\[
\int_0^x f = n \int_0^T f + \int_{nT}^{nT + r} f,
\qquad
\Bigl| \int_{nT}^{nT+r} f \Bigr| \leq T \sup_{\intcc{0}{T}} \abs f
= C .
\]
Then $\frac 1x \int_0^x f = \frac{nT}{x}\cdot\frac 1T \int_0^T f +
O\bigl(\frac 1x\bigr)$, and $\frac{nT}{x} \to 1$: the limit is
$\frac1T \int_0^T f$.
\end{solution}

\begin{solution}{exo:b1:integration:7}
Let $g(x) = f(x) - x$: continuous, with
\[
\int_0^1 g = \int_0^1 f - \frac12 = 0 .
\]
If $g$ never vanished, the intermediate value theorem would force a
constant sign (a continuous function on an interval taking both signs
vanishes); say $g > 0$. Then, by strict positivity
(\cref{thm:b1:integration:props}~(4) applied to $g > 0$, giving
$\int g > 0$): contradiction with $\int g = 0$. So $g(c) = 0$ for
some $c$: $f(c) = c$.
\end{solution}

\begin{solution}{exo:b1:integration:8}
\begin{enumerate}
  \item By parts with $u' = \sin t$, $v = \sin^{n-1} t$:
        \[
        W_n = \bigl[-\cos t\sin^{n-1}t\bigr]_0^{\pi/2}
        + (n-1)\int_0^{\pi/2} \cos^2 t\,\sin^{n-2} t\,\dd t
        = (n-1)(W_{n-2} - W_n),
        \]
        so $nW_n = (n-1)W_{n-2}$. From $W_0 = \frac\pi2$, $W_1 = 1$:
        \[
        W_{2p} = \frac{(2p-1)(2p-3)\cdots 1}{(2p)(2p-2)\cdots 2}\,
        \frac{\pi}{2} = \frac{(2p)!}{4^p (p!)^2}\,\frac\pi2,
        \qquad
        W_{2p+1} = \frac{(2p)(2p-2)\cdots 2}{(2p+1)(2p-1)\cdots 3}
        = \frac{4^p (p!)^2}{(2p+1)!} .
        \]
  \item On $\intoo{0}{\frac\pi2}$, $0 < \sin t < 1$, so $\sin^{n+1} <
        \sin^n$ and $(W_n)$ is (strictly) decreasing, positive.
        Sandwiching with the recurrence:
        \[
        \frac{n}{n+1} = \frac{W_{n+1}}{W_{n-1}}
        \leq \frac{W_{n+1}}{W_n} \leq 1
        \quad\implies\quad \frac{W_{n+1}}{W_n} \to 1 .
        \]
        Invariant: $a_n = (n+1)W_{n+1}W_n$ satisfies $a_n = a_{n-1}$
        by the recurrence $(n+1)W_{n+1} = nW_{n-1}$, so $a_n = a_0 =
        1 \cdot W_1 W_0 = \frac\pi2$. Then
        \[
        n W_n^2 \sim (n+1) W_{n+1} W_n = \frac\pi2
        \quad\implies\quad
        W_n \sim \sqrt{\frac{\pi}{2n}} .
        \]
\end{enumerate}
\end{solution}

\begin{solution}{exo:b1:integration:9}
\begin{enumerate}
  \item On $\intoo{0}{\pi}$: $x > 0$, $a - bx = b(\frac ab - x) =
        b(\pi - x) > 0$ and $\sin x > 0$, so the integrand is $> 0$
        and $I_n > 0$ (strict positivity). Bound: on
        $\intcc{0}{\pi}$, $x \leq \pi$ and $a - bx \leq a$, so $P
        \leq \frac{\pi^n a^n}{n!}$ and $I_n \leq
        \pi\,\frac{(\pi a)^n}{n!}$, which tends to $0$ (the
        factorial beats the geometric term: it is the general term
        of the convergent exponential series, cf.\
        \cref{ex:b1:seq:e}); in particular $I_n < 1$ for large $n$.
  \item Expand $x^n(a - bx)^n = \sum_{j=0}^{n} \binom nj a^{n-j}
        (-b)^j x^{n+j}$: so $P = \frac{1}{n!}\sum_j c_j x^{n+j}$
        with integer $c_j$. Then $P^{(k)}(0) = 0$ for $k < n$
        (valuation) and, for $n \leq k \leq 2n$, $P^{(k)}(0) =
        \frac{k!}{n!} c_{k-n}$, an integer since $n! \mid k!$.
        Moreover $P(\pi - x) = P(x)$ (substitute: $\pi - x$ swaps
        the factors, using $a - b(\pi - x) = bx$), so $P^{(k)}(\pi)
        = \pm P^{(k)}(0)$: integers as well.
  \item With $Q = P - P'' + P^{(4)} - \dots$ (finite: $P$ has degree
        $2n$): $Q + Q'' = P$, and
        \[
        \bigl(Q'\sin x - Q\cos x\bigr)' = (Q + Q'')\sin x = P\sin x .
        \]
        Hence $I_n = \bigl[Q'\sin x - Q\cos x\bigr]_0^{\pi} = Q(\pi)
        + Q(0)$, a sum of values $P^{(2k)}$ at $0$ and $\pi$: an
        integer by (2).
  \item For $n$ large, $I_n$ is an integer with $0 < I_n < 1$:
        impossible. The assumption $\pi = \frac ab$ fails: $\pi$ is
        irrational.
\end{enumerate}
\end{solution}

\begin{solution}{exo:b1:integration:10}
$C^1$ case: by parts,
\[
\int_a^b f(t)\sin\lambda t\,\dd t
= \Bigl[-f(t)\frac{\cos\lambda t}{\lambda}\Bigr]_a^b
+ \frac{1}{\lambda}\int_a^b f'(t)\cos\lambda t\,\dd t,
\]
bounded in absolute value by $\frac{2\sup\abs f + (b - a)\sup\abs{f'}}
{\lambda} \to 0$.

Continuous case: let $\varepsilon > 0$ and pick a step function
$\varphi$ with $\abs{f - \varphi} \leq \varepsilon$
(\cref{thm:b1:integration:approx} provides $\varphi \leq f \leq
\psi$ with gap $\leq\varepsilon$; take $\varphi$). Then
\[
\Bigl|\int_a^b f\sin\lambda t\Bigr|
\leq \int_a^b \abs{f - \varphi}
+ \Bigl|\int_a^b \varphi \sin\lambda t\Bigr|
\leq (b-a)\varepsilon
+ \sum_i \abs{c_i}\,\Bigl|\int_{x_{i-1}}^{x_i} \sin\lambda t\,\dd
t\Bigr| ,
\]
and each $\bigl|\int \sin \lambda t\bigr| = \bigl|\frac{\cos\lambda
x_{i-1} - \cos\lambda x_i}{\lambda}\bigr| \leq \frac{2}{\lambda}$:
the second term tends to $0$. Hence the limsup is $\leq
(b-a)\varepsilon$ for every $\varepsilon$: the limit is $0$.
\end{solution}
